% /Users/pikachu/books/payments-book/src/parts/part4/ch22-reconciliation.tex
\chapter{Сверка: где теряются деньги}
\label{ch:reconciliation}
\margintoc

\begin{kaobox}[title=Что вы узнаете из этой главы]
  Сверка нужна не затем, чтобы свести красивые отчёты, а затем, чтобы
  деньги не исчезали между авторизацией, клирингом и расчётом. В этой
  главе мы разберём три разрыва, на которых обычно распадается история
  одной операции; посмотрим, как читать клиринговые файлы Visa и
  Mastercard; отличим нетто-расчёт от валового; обсудим исключения,
  поздние представления и мультивалютные случаи; а затем поймём,
  почему автоматизация сверки начинается не с дашборда, а с
  дисциплины сопоставления.
\end{kaobox}

Деньги почти никогда не \enquote{теряются} в одной точке.
Обычно они растворяются в переходах между системами, каждая из которых
честно хранит свою версию события. В журнале авторизаций покупка уже
разрешена, в клиринговом файле она ещё не появилась, а в банковской
выписке от неё осталась только нетто-сумма после комиссий и взаимных
вычетов. Задача сверки в том и состоит, чтобы из этих разрозненных
следов снова собрать одну операцию.

Внутренний реестр обязательств из главы \ref{ch:internal-ledger}
отвечает на вопрос, какое денежное состояние видит сам продукт.
Сверка задаёт следующий вопрос: совпадает ли эта внутренняя картина
с внешними файлами схемы, эквайера и банка. Если перепутать эти два
уровня, команда начнёт сравнивать несопоставимые величины и искать
потери там, где есть лишь разница между моделями учёта.

\section{Три разрыва}
\label{sec:recon-gaps}

Почему одна и та же покупка выглядит по-разному в разных системах?
Потому что карточная операция живёт не одним событием, а цепочкой.
Авторизация фиксирует право провести операцию, клиринг закрепляет
её окончательный денежный облик, а расчёт показывает уже фактическое
движение средств с комиссиями и взаимными вычетами. Каждая стадия
оставляет собственный след, и сверка начинается с признания того, что
эти следы по определению не обязаны быть идентичными.

Первый разрыв возникает между авторизацией и клирингом. Ресторан может
авторизовать счёт на одну сумму, а после чаевых отправить в клиринг
другую; отель сначала блокирует предварительный депозит, а потом
предъявляет итоговую стоимость проживания; авиакомпания доначисляет
часть сборов уже после исходной покупки. Поэтому сверка на этом уровне
начинается не с требования буквального равенства, а с поиска правильной
пары: для каждой клиринговой записи надо найти соответствующую
авторизацию и проверить, укладывается ли расхождение в допустимые
пределы \sidecite{visa-core-rules}.

Второй разрыв возникает между клирингом и расчётной позицией. За день
может пройти сто транзакций, но в расчётном отчёте эквайер увидит уже
не их полный перечень, а свёрнутую итоговую сумму. Именно здесь многие
команды впервые чувствуют, что \enquote{деньги не сходятся}, хотя на
самом деле они сравнивают поток операций с результатом его
нетто-свёртки. Схема объединяет покупки, возвраты, комиссии и другие
корректировки в одну расчётную величину, поэтому сумма клиринга и сумма
расчёта почти никогда не совпадают буквально.

Третий разрыв возникает между расчётной позицией и банковской выпиской.
Число в системе схемы может уже считаться закрытым, а движение по
корреспондентскому счёту пройти позже, в другом окне дня или даже в
иной валюте после пересчёта. Здесь сверка работает уже не с
транзакциями как таковыми, а с сопоставлением расчётного файла и
банковской выписки по счёту урегулирования. Поэтому прежде чем искать
потери, полезно точно назвать слой, на котором вы их ищете. Иначе спор
пойдёт не о деньгах, а о разных способах описать один и тот же поток.

\section{TC33 и IPM}
\label{sec:recon-tc33-ipm}

Как читать клиринговые файлы схем? Прежде всего не принимать их за
ещё один бухгалтерский отчёт. Это рабочие документы расчётного контура,
и читать их нужно как машинный след завершённых транзакций. Форматы у
Visa и Mastercard устроены по-разному, хотя задача у них одна:
передать операцию из клиринга в расчёт.

Файл Transaction Clearing 33 (TC33)\sidenote{TC33 — формат клирингового
файла Visa, в котором эквайер передаёт завершённые транзакции в базовый
расчётный контур.} используется Visa для передачи клиринговых данных
\sidecite{visa-core-rules}. Структура файла иерархична: заголовок,
группы транзакций, сами транзакции, сводные записи и хвост с
контрольными суммами. Для сверки важны не столько названия блоков,
сколько две вещи: контрольные суммы, которые позволяют сразу увидеть
повреждённый файл, и поля, по которым клиринг можно связать с
предыдущей авторизацией.

В TC33 особенно полезны идентификаторы, по которым можно связать
клиринг с авторизацией: сумма, дата, STAN, Terminal ID, Merchant ID,
код ответа авторизации и признак схемной тарифной программы. Опасность здесь
в том, что дата транзакции и дата обработки не обязаны совпадать.
Поэтому сопоставление нельзя строить только по календарю; ему нужен
набор технических признаков, который переживает сдвиг операционного дня.

Mastercard использует формат Integrated Product Messages (IPM)\sidenote{IPM — клиринговый формат Mastercard для передачи транзакций в расчётный контур.}.
Он решает ту же задачу, но в собственной структуре со служебными
блоками и дополнительными полями, которые Mastercard добавляет к
стандартным карточным данным \sidecite{mc-rules}. Внутри IPM особую
роль играют блоки частных данных: именно в них живёт специфика
конкретной программы, которую нельзя выразить одним только базовым
набором полей.

На практике это означает простую, но важную вещь: парсер должен
собрать транзакцию, извлечь ключевые поля и передать их в механизм
сопоставления, не делая вид, будто TC33 и IPM — это один и тот же
формат. Для каждой схемы нужен свой адаптер или хотя бы аккуратный
промежуточный слой. Иначе ошибка возникает ещё до бухгалтерии, на
этапе разбора файла.

Именно поэтому клиринговый файл не равен отчёту для бухгалтера. Это
сырьё для учётной картины, а не её готовая и человекочитаемая версия.

\section{Нетто- и валовой расчёт}
\label{sec:recon-net-vs-gross}

Почему сумма транзакций почти никогда не совпадает с тем, что пришло на
расчётный счёт? Потому что карточные системы обычно показывают не вал
операций, а результат их свёртки. По умолчанию здесь работает
нетто-расчёт (net settlement), а валовой расчёт (gross settlement)
остаётся скорее исключением.

В модели нетто-расчёта схема сводит все позиции дня в одну итоговую
сумму. За множеством покупок, возвратов и комиссий остаётся одна
очищенная позиция, которую и видит эквайер. Такой подход удобен
операционно: он уменьшает число отдельных переводов и делает расчётный
контур дешевле и спокойнее.

Валовой расчёт устроен иначе: каждая операция или небольшая группа
операций идёт отдельно. Такой режим прозрачнее, потому что деньги не
растворяются в общей дневной сумме, но расплачивается он другой ценой:
учёт становится более детальным, а дисциплина сверки — более строгой.

Отсюда следует практическое правило. Если схема работает в нетто-режиме,
бессмысленно искать в банковской выписке сумму всех покупок за день.
Сначала сверяют клиринговый поток с расчётной позицией, и только потом
расчётную позицию — с банковской выпиской.

\begin{kaobox}[title=На практике]
  Типовая ошибка при первичной настройке сверки состоит в том, что
  авторизационный журнал сравнивают сразу с банковской выпиской. Это
  неправильно: авторизация ещё не стала клирингом, банковская
  выписка уже отражает нетто-расчёт. Правильная цепочка сравнения
  выглядит так: авторизация против клиринга, клиринг против расчётной
  позиции, расчётная позиция против банковской выписки.
\end{kaobox}

\section{Исключения}
\label{sec:recon-exceptions}

Что делать с транзакциями, которые не сходятся автоматически? Не
прятать их в общей массе и не маскировать под норму. Зрелая
сверка начинается с явной классификации: каждое несоответствие
получает понятный тип и свой маршрут обработки — либо автоматический,
либо ручной.

Чаще всего первым всплывает расхождение суммы. В ресторане
авторизовали одну сумму, а в клиринге пришла другая; в гостинице
предварительная блокировка не совпала с окончательным чеком. Такая
запись не требует паники, но требует проверки допустимого отклонения.
Если порог превышен, считать её штатной уже нельзя \sidecite{visa-core-rules}.

Другая группа исключений — клиринг без авторизации. Такое бывает при
принудительном представлении, техническом сбое или в редких сценариях,
когда авторизация прошла вне основного контура записи. Здесь важнее
всего не автоматический ярлык, а объяснение происхождения: перед нами
допустимый сценарий, ошибка маршрутизации, технический хвост или
признак риска.

Обратная проблема — авторизация без клиринга. Операция была разрешена,
но в клиринг так и не ушла. Это может означать ошибку маршрутизации,
сбой у мерчанта или истечение окна представления. Если такую запись
вовремя не закрыть, держатель карты увидит только авторизационную
блокировку, а мерчант так и не получит выручку.

Наконец, есть позднее представление в клиринг: транзакция дошла до
схемы позже, чем позволяют обычные сроки. Реакция зависит от правил
сети: запись могут принять, отклонить или провести на иных
экономических условиях. Для команды сверки это означает, что одной
бухгалтерской классификации недостаточно: нужно ещё понять, нарушено
ли окно представления по правилам схемы \sidecite{visa-dispute-rules}.

\section{DCC и мультивалютный расчёт}
\label{sec:recon-dcc}

Как сверка меняется, если в операции участвуют разные валюты? На первый
взгляд DCC меняет только сумму на экране покупателя, но для учёта
меняется сама точка, в которой фиксируется курс. В сценарии
динамической конвертации валюты (Dynamic Currency Conversion, DCC)
клиент видит сумму в своей валюте ещё на кассе, а мерчант зарабатывает
на курсовой марже.

Для расчёта это простая только снаружи история. Если конвертация уже
произошла на кассе, повторно пересчитывать сумму нельзя. В клиринге и
расчётной позиции должны быть корректно отражены и валюта транзакции,
и валюта урегулирования, иначе система начнёт сравнивать не те величины.

Для мерчанта такой поток несёт ещё и валютный риск. Если расчёт идёт не
в той же валюте, в какой прошла покупка, разница между курсом
транзакции и курсом расчёта попадает прямо в финансовый результат.
Поэтому сверка должна уметь отделять курсовое отклонение от настоящей
ошибки учёта.

Именно здесь становится видно, что DCC и мультивалютный расчёт —
не декоративная надстройка, а отдельный слой правил. Без него сверка
быстро превращается в поиск фантомных расхождений.

\section{Автоматизация сверки}
\label{sec:recon-automation}

Как перестать сводить файлы вручную в таблицах? Нужен механизм
сопоставления\sidenote{Механизм сопоставления — компонент сверки,
который сначала ищет точные совпадения, затем допускает частичные и
отправляет остаток в обработку исключений.}, который умеет находить
одну и ту же операцию в разных представлениях и не теряет осторожности
там, где совпадение неочевидно. В простейшем виде такой компонент
сначала ищет точные совпадения, затем допускает более мягкие, а
остаток отправляет в очередь исключений.

Обычно он работает в несколько проходов: сначала строгий матч по
идентификаторам и сумме, затем более мягкая проверка по дате и другим
полям, затем очередь исключений. Ограничение здесь не техническое, а
организационное: чем слабее критерии сопоставления, тем аккуратнее
должна быть ручная верификация.

Готовый сервис сверки имеет смысл там, где нет прямого доступа к
схемам или объём транзакций не оправдывает собственной разработки.
Свой механизм обычно строят там, где есть сложный мультивалютный
контур, собственная ERP и специфические правила нетто-расчёта.

Отдельная часть этой темы — наблюдаемость. Если команда не смотрит на
долю сопоставленных записей, время закрытия дня и размер очереди
исключений, она не понимает, работает ли система как надо или просто
молча накапливает проблемы.

\begin{kaobox}[title=На практике]
  Самые упрямые ложные расхождения в сверке связаны с датой отсечки.
  Схемный расчёт живёт по операционному времени закрытия дня, а не по
  календарной дате. Поэтому транзакция, прошедшая поздно вечером,
  может попасть в следующий расчётный день. Если система жёстко
  привязана к календарной дате, фиктивные расхождения будут появляться
  регулярно, особенно в конце недели и перед праздниками.
\end{kaobox}

\begin{chaptersummary}
\begin{itemize}
  \item Сверка закрывает три разрыва: авторизация против клиринга,
        клиринг против расчётной позиции, расчётная позиция против
        банковской выписки.
  \item TC33 и IPM передают клиринговые данные в разных форматах, но
        требуют одной дисциплины: разобрать файл, извлечь ключевые поля
        и связать их с авторизацией.
  \item Нетто-расчёт даёт одну итоговую позицию за период, поэтому
        сумма транзакций не равна банковской выписке.
  \item Основные исключения сводятся к расхождению суммы, клирингу без
        авторизации, авторизации без клиринга и позднему представлению.
  \item DCC и мультивалютный расчёт требуют отдельного учёта валютного
        риска и запрета на повторную конвертацию.
  \item Автоматизация сверки начинается с механизма сопоставления,
        очереди исключений и наблюдаемости.
\end{itemize}
\end{chaptersummary}
